% ---------------------------------------------------------------------------
% paper-preamble.tex -- shared by paper.tex and all section files
% Compiles with pdfLaTeX, XeLaTeX and LuaLaTeX.
% ---------------------------------------------------------------------------
\usepackage{iftex}
\ifPDFTeX
  \usepackage[T1]{fontenc}
  \usepackage[utf8]{inputenc}
  \usepackage{lmodern}
\else
  \usepackage[no-math]{fontspec}
  \IfFontExistsTF{DejaVuSansMono.ttf}{%
    \setmonofont{DejaVuSansMono.ttf}[Scale=MatchLowercase,
      BoldFont=DejaVuSansMono-Bold.ttf, ItalicFont=DejaVuSansMono-Oblique.ttf,
      BoldItalicFont=DejaVuSansMono-BoldOblique.ttf]%
  }{}
\fi
\usepackage[a4paper,margin=2.7cm]{geometry}
\usepackage{amsmath,amssymb,amsthm,mathtools}
\usepackage{enumitem}
\usepackage{booktabs,longtable,array}
\usepackage{xcolor}
\usepackage{url}
\usepackage[hidelinks]{hyperref}
\usepackage[capitalise,nameinlink,noabbrev]{cleveref}

% --- theorem-like environments (one counter, numbered by section) -----------
\theoremstyle{plain}
\newtheorem{theorem}{Theorem}[section]
\newtheorem{proposition}[theorem]{Proposition}
\newtheorem{lemma}[theorem]{Lemma}
\newtheorem{corollary}[theorem]{Corollary}
\newtheorem{claim}[theorem]{Claim}
\newtheorem{littheorem}[theorem]{Theorem}
\newtheorem*{theoremA}{Theorem A}
\newtheorem*{theoremB}{Theorem B}
\newtheorem*{corollaryC}{Corollary C}
\theoremstyle{definition}
\newtheorem{definition}[theorem]{Definition}
\newtheorem{hypothesis}[theorem]{Hypothesis}
\newtheorem{example}[theorem]{Example}
\theoremstyle{remark}
\newtheorem{remark}[theorem]{Remark}
\crefname{littheorem}{Theorem}{Theorems}
\Crefname{littheorem}{Theorem}{Theorems}
\crefname{hypothesis}{Hypothesis}{Hypotheses}
\Crefname{hypothesis}{Hypothesis}{Hypotheses}
\crefname{claim}{Claim}{Claims}
\Crefname{claim}{Claim}{Claims}
\numberwithin{equation}{section}
\setlength{\emergencystretch}{3em}

% --- references to the formal development ------------------------------------
% \formal{Full.Lean.Name}: the Lean declaration that states (and proves) a result.
% A ragged line break is allowed before the tag, and the name may break after dots and underscores.
\DeclareUrlCommand\formalname{\def\UrlFont{\ttfamily}\def\UrlBreaks{\do\.\do\_}\def\UrlBigBreaks{}}
\newcommand{\formal}[1]{\ifhmode\hfil\penalty50\hfilneg\fi{\footnotesize\textsf{[}\formalname{#1}\textsf{]}}}
% \Fmark: a result of Part I whose proof is only outlined there and given in Part II (see Section 1.4)
\newcommand{\Fmark}{\textsuperscript{\textsf{F}}}
% \formallabel{Full.Lean.Name}: label lean:<name> for cross-references in Part II, and print the tag
\newcommand{\formallabel}[1]{\label{lean:#1}\formal{#1}}
% \lref{Full.Lean.Name}: the result of this paper with label lean:<name> if there is one; otherwise the
% bracketed name of the formal result (whose proof is then only contained in the formal development)
\makeatletter
\DeclareRobustCommand{\lref}[1]{\@ifundefined{r@lean:#1}{\formal{#1}}{\cref{lean:#1}}}
\DeclareRobustCommand{\Lref}[1]{\@ifundefined{r@lean:#1}{\formal{#1}}{\Cref{lean:#1}}}
\makeatother
\pdfstringdefDisableCommands{\def\lref#1{#1}\def\Lref#1{#1}\def\formal#1{#1}}
\newenvironment{proofoutline}{\begin{proof}[Outline of proof]}{\end{proof}}

% --- notation -----------------------------------------------------------------
\newcommand{\N}{\mathbb{N}}
\newcommand{\Z}{\mathbb{Z}}
\newcommand{\Q}{\mathbb{Q}}
\newcommand{\R}{\mathbb{R}}
\newcommand{\C}{\mathbb{C}}
\newcommand{\F}{\mathbb{F}}
\newcommand{\A}{\mathbb{A}}
\newcommand{\Qbar}{\overline{\mathbb{Q}}}
\newcommand{\Ql}{\mathbb{Q}_\ell}
\newcommand{\Zl}{\mathbb{Z}_\ell}
\newcommand{\Qlbar}{\overline{\mathbb{Q}}_\ell}
\newcommand{\OO}{\mathcal{O}}
\newcommand{\Sha}{\mathrm{III}}
\renewcommand{\Re}{\operatorname{Re}}
\renewcommand{\Im}{\operatorname{Im}}
\DeclareMathOperator{\GL}{GL}
\DeclareMathOperator{\SL}{SL}
\DeclareMathOperator{\PGL}{PGL}
\DeclareMathOperator{\Sp}{Sp}
\DeclareMathOperator{\GSp}{GSp}
\DeclareMathOperator{\Uni}{U}
\DeclareMathOperator{\GU}{GU}
\DeclareMathOperator{\SO}{SO}
\DeclareMathOperator{\Ort}{O}
\DeclareMathOperator{\Gal}{Gal}
\DeclareMathOperator{\Hom}{Hom}
\DeclareMathOperator{\End}{End}
\DeclareMathOperator{\Aut}{Aut}
\DeclareMathOperator{\Ind}{Ind}
\DeclareMathOperator{\Res}{Res}
\DeclareMathOperator{\Frob}{Frob}
\DeclareMathOperator{\Fr}{Fr}
\DeclareMathOperator{\tr}{tr}
\DeclareMathOperator{\Tr}{Tr}
\DeclareMathOperator{\Nm}{N}
\DeclareMathOperator{\vol}{vol}
\DeclareMathOperator{\covol}{covol}
\DeclareMathOperator{\rank}{rank}
\DeclareMathOperator{\ord}{ord}
\DeclareMathOperator{\Spec}{Spec}
\DeclareMathOperator{\Pic}{Pic}
\DeclareMathOperator{\Lie}{Lie}
\DeclareMathOperator{\cond}{cond}
\DeclareMathOperator{\supp}{supp}
\DeclareMathOperator{\lcm}{lcm}
\DeclareMathOperator{\len}{length}
\newcommand{\one}{\mathbf{1}}
\newcommand{\eps}{\varepsilon}
\newcommand{\sfname}[1]{\mathsf{#1}}   % names of propositions of the development, e.g. \sfname{VWin}
